\documentclass[spanish,letterpaper,oneside]{article}

\def\documenttitle {Trabajo Final Integrador: Propuesta de Intervención Pedagógica Transformadora}
\def\documentsubtitle {Presentación Ejecutiva · Cohorte SD26}
\def\documentsubject {Actividad 8 · Síntesis Visual}
\def\documentauthor {Mtro. Martín Jonathan de la Cruz Muñoz}
\def\coursename {Fundamentos para la enseñanza y el aprendizaje I}
\def\coursecode {MAEM-SD26}

\def\universityname {Instituto de Investigación, Innovación y Estudios de Posgrado para la Educación}
\def\universityfaculty {Maestría en Aprendizaje y Enseñanza de las Matemáticas}
\def\universitydepartment {Fundamentos para la enseñanza y el aprendizaje I}
\def\universitydepartmentimage {img/iiiepe/marca-iiiepe.png}
\def\universitydepartmentimagecfg {height=1.75cm}
\def\universitylocation {Monterrey, Nuevo León}
\def\studentfullname {Mtro. Martín Jonathan de la Cruz Muñoz}
\def\studentid {00874}
\def\studentemail {martin.delacruz@campus.iiiepe.edu.mx}
\def\programperiod {2do periodo, septiembre--diciembre de 2026 (SD26)}
\def\programmodality {Modalidad presencial}
\def\programshort {MAEM}

\def\authortable {
  \begin{tabular}{ll}
    \textbf{Alumno:} & \studentfullname \\
    Matrícula: & \studentid \\
    Unidad: & \coursename \\
    Periodo: & \programperiod
  \end{tabular}
}

\input{template}

\usepackage{helvet}
\renewcommand{\familydefault}{\sfdefault}
\usepackage{tikz,longtable,array}
\usetikzlibrary{arrows.meta,positioning}
\definecolor{iiiepeNavy}{HTML}{163B5C}
\definecolor{iiiepeBlue}{HTML}{1F6F9F}
\definecolor{iiiepePale}{HTML}{EEF5F8}

\begin{document}\templatePortrait\templatePagecfg\templateFinalcfg\setstretch{1.25}

\section*{Síntesis Ejecutiva · Actividad 8}
\textbf{Trabajo Final Integrador: Propuesta de Intervención Pedagógica Transformadora}

\subsection*{1. Diagnóstico y Planteamiento Central}
La presente exposición aborda la articulación de los marcos pedagógicos contemporáneos en el aula de matemáticas:
\begin{itemize}
  \item Eje temático: Integración global de fundamentos epistemológicos, identidad docente, aprendizaje profundo y modelo bioecológico en una secuencia didáctica..
  \item Planteamiento central: La enseñanza matemática trasciende el adiestramiento mecánico cuando incorpora la indagación y la justificación argumentada \cite{delors1996educacion,scott2015futuro}.
\end{itemize}

\subsection*{2. Dimensiones y Artefacto de Aplicación}
{\small\renewcommand{\arraystretch}{1.25}
\begin{longtable}{|p{0.35\linewidth}|p{0.6\linewidth}|}\hline
\rowcolor{iiiepePale} \textbf{Dimensión} & \textbf{Acción Didáctica en el Aula} \\\hline
Fundamentación Teórica & Vinculación conceptual con las ciencias del aprendizaje \cite{sawyer2008optimizing}. \\\hline
Mediación Matemática & Diseño de tareas abiertas con representaciones múltiples. \\\hline
Impacto Profesional & Transferencia reflexiva al contexto laboral y comunitario. \\\hline
\end{longtable}}

\subsection*{3. Conclusión y Veredicto Pedagógico}
En conclusión, la práctica docente se consolida al articular rigor disciplinar y enfoque humanista. Por tanto, esta propuesta formativa promueve la autonomía y equidad educativa en los estudiantes.

\end{document}
